\section{量子アルゴリズムにおけるサブルーチン}
\label{sec:role-of-quantum-subroutines}

\subsection{サブルーチンの役割}
\label{subsec:quantum-subroutine-role}

サブルーチンとは，特定の変換を実行し，複数のアルゴリズムから繰り返し利用される処理単位である．
古典アルゴリズムでは，サブルーチンの入力と出力は通常，ビット列や数値として与えられる．
これに対し，量子サブルーチンでは，入力が重ね合わせ状態であっても線形性を保ったまま処理する必要がある．

例えば，あるユニタリ演算子$U$の固有状態$\ket{u_j}$に対応する位相$\phi_j$を求める処理を考える．
入力が一つの固有状態であれば
\begin{align*}
  \ket{0}^{\otimes t}\ket{u_j}
  \longmapsto
  \ket{\widetilde{\phi}_j}\ket{u_j}
\end{align*}
と変換すればよい．
しかし，入力が固有状態の重ね合わせ
\begin{align*}
  \ket{\psi}
  \coloneq
  \sum_j c_j\ket{u_j}
\end{align*}
である場合，量子サブルーチンは
\begin{align*}
  \ket{0}^{\otimes t}\ket{\psi}
  \longmapsto
  \sum_j c_j\ket{\widetilde{\phi}_j}\ket{u_j}
\end{align*}
と作用する．
これにより，固有状態を測定して一つに確定させることなく，
各固有状態と対応する位相の関係を量子状態の中に保持できる．

\subsection{補助量子ビットと逆演算}
\label{subsec:ancilla-and-uncomputation}

量子サブルーチンでは，途中の計算結果を保持するために補助量子ビットを用いる．
補助量子ビットが不要になった後も計算結果が残っていると，
目的のレジスタと量子もつれを持ち，その後の干渉を妨げることがある．
そこで，必要な情報を別のレジスタへ反映した後にサブルーチンの逆演算を適用し，
補助量子ビットを初期状態へ戻す．
この処理を\textbf{計算の巻き戻し}または\textbf{アンコンピュテーション}という．

HHLアルゴリズムでは，量子位相推定によって固有値の情報を位相レジスタへ書き込み，
その値に応じた制御回転を補助量子ビットへ作用させる．
その後，逆量子位相推定によって位相レジスタを$\ket{0}^{\otimes t}$へ戻す．
このように，サブルーチンは単に値を計算するだけでなく，
重ね合わせ，量子もつれ，逆演算を保ったまま他の処理へ接続できることが重要である．

\subsection{本章で扱う処理}
\label{subsec:subroutines-in-this-chapter}

本章では，次の三つの処理を扱う．

\begin{enumerate}[(1)]
  \item 位相を計算基底のビット列へ変換する\textbf{量子フーリエ変換}
  \item ユニタリ演算子の固有位相を推定する\textbf{量子位相推定}
  \item 目的とする測定結果の確率振幅を大きくする\textbf{振幅増幅}
\end{enumerate}

量子フーリエ変換は量子位相推定の内部で用いられる．
量子位相推定と振幅増幅は独立した処理であるが，
両者を組み合わせると，成功確率そのものを推定する振幅推定も構成できる．
さらに，量子位相推定と振幅増幅はHHLアルゴリズムの固有値処理と成功確率の改善に利用される．
